\documentclass[12pt,a4paper]{ctexart}
\usepackage{amsmath,amssymb}
\usepackage{booktabs}
\usepackage{array}
\usepackage{geometry}
\usepackage{enumitem}
\usepackage[hidelinks]{hyperref}
\geometry{margin=2.5cm}
\setlength{\parskip}{0.4em}
\setlength{\parindent}{2em}
\linespread{1.25}
\renewcommand{\arraystretch}{1.15}
\setcounter{secnumdepth}{-1}
\ctexset{
  section = {format = {\heiti \zihao{3}}},
  subsection = {format = {\heiti \zihao{4}}},
  subsubsection = {format = {\heiti \zihao{-4}}},
}

\begin{document}

\begin{center}
{\heiti\zihao{-2} A 题论文评审报告\par}
\vspace{0.4em}
{\heiti\zihao{4} （国赛评委视角 · 内部评审稿）\par}
\vspace{0.3em}
{\zihao{5} 评审对象：\texttt{A题论文.tex}（PDF 共 18 页，摘要页 1 页，正文约 16 页）\quad 日期：2026-09-12\par}
\end{center}

\vspace{0.6em}

\section{一、总体评价}

这是一篇\textbf{具备省一等奖以上竞争力}的论文。模型与题设严格对应、数值格式严谨、检验体系超出多数参赛队水平。评审期间已完成三项独立核查：论文引用的附件数据（附件 1 共 241 点、附件 2 共 145 点、末值 $50.165\,^\circ$C、67 h 起半径恒为 1.198 cm、收缩 6.4\%/40.1\%）全部准确；表 1--6 与 result1--4.xlsx 逐格吻合；摘要中的关键数字与表格一致，7.2 节 Richardson 外推（$\approx57.05$ h）、7.6 节 Arrhenius 比值（0.72）、$\mathrm{Bi}=1.389$ 等推导复查无误。\textbf{未发现数学或数据硬伤}。

主要失分点集中在\textbf{表述的清晰程度}：全文无图（已按团队要求暂缓处理）、缺灵敏度分析、摘要已到页满且缺问题 2 结果、7.6 节「反证」措辞存在逻辑风险，另有若干公式记号与验证完备性问题。逐项见下文。

\section{二、评审依据}

\begin{enumerate}[label=(\arabic*),leftmargin=3em]
\item \textbf{《全国大学生数学建模竞赛章程》四项评奖标准}：假设的合理性、建模的创造性、结果的正确性、文字表述的清晰程度；
\item \textbf{《全国大学生数学建模竞赛论文格式规范（2026 年修订稿）》}（见文末来源）；
\item 独立核查：附件数据逐格核对、公式推导复查、PDF 版式与结果文件结构检查，并交叉引用了本目录下《问题四分析\_发现的问题.md》的代码审查结论。
\end{enumerate}

\section{三、问题清单（汇总）}

\begin{center}
{\zihao{5}\begin{tabular}{clcl}
\toprule
编号 & 问题 & 严重度 & 位置 \\
\midrule
1 & 全文无图，18 页纯表（暂缓，待配图方案确定） & 高（暂缓） & 全文 \\
2 & 缺灵敏度分析：恒温段假设（$T_a,C_a$ 取附件 1 末值恒定）及 $h,h_m$ 对 $t^*$ 的影响未量化 & 高 & 7.6 之后 \\
3 & 摘要已到页满且缺问题 2 的结果 & 高 & 摘要页 \\
4 & 问题 4 空间误差 $\approx1.8\times10^{-4}$，超出四位小数分辨率，表 6 前几行第 4 位小数存疑 & 中 & 7.2 / 表 6 \\
5 & 问题 4 的 $t^*$ 未报网格收敛（问题 3 已报） & 中 & 7.2 \\
6 & 「收缩使烘干缩短 6.3 h」归因混杂：物性差异与几何收缩两个因素未分离 & 中 & 6.4 \\
7 & 式 (cn) 与式 (L) 的 $v_i$ 记号重复，量纲不一致（代码正确，属公式笔误） & 中 & 5.4 \\
8 & 7.6 节「反证」措辞有循环论证风险 & 中 & 7.6 \\
9 & $h,h_m$ 沿用附录 2 取值未在 5.2 节显式声明 & 中 & 5.2 / H4 \\
10 & 一级标题/题目字体与通行惯例不符 & 低 & 全文格式 \\
11 & 表 6 距离列的欧拉视角解释未明示 & 低 & 5.3 \\
12 & result3/4 末行 $t^*$ 与 60 s 间隔序列的关系未说明 & 低 & 附录 5 \\
13 & $\delta\sim D/h_m$ 与 $\sqrt{Dt}$ 两个尺度并存未解释 & 低 & 5.5 / 6.1 \\
14 & Picard 两轮迭代的耦合残差无量化 & 低 & 5.4 \\
15 & 文末无「结论」清单（四问答案汇总） & 低 & 文末 \\
16 & 8.1「模型优点」未突出本文三个特色点 & 低 & 8.1 \\
\bottomrule
\end{tabular}}
\end{center}

\section{四、逐项评审}

\subsection{4.1 假设的合理性（优）}

H1--H8 每条都有依据：长径比 12.5 与 $(L/R)^2\approx156$ 支撑一维化；$h,h_m$ 常数有文献支撑；H6 用附件 2 数据量化收缩影响。符合「必要假设、切合题意」的要求。

\textbf{H5（潜热有效参数观点）是全文最大胆的假设}，但 7.6 节做了数量级论证且 H5 处显式交叉引用，处理得当。唯一疏漏（\textbf{问题 9}）：问题 2--4 题文未另行给出 $h,h_m$，论文沿用了附录 2 的值，评委能推断出来，但建议在 5.2 节加一句「题文未另给对流系数，沿用附录 2 取值」。

\subsection{4.2 建模的创造性（良）}

本题参数已由题目锁定，发挥空间有限，论文的加分点选得准：守恒型离散+表面节点置于真实表面（守恒闭合到 $10^{-18}$）、问题 4 随体坐标处理（干基含水率定义于固相骨架 $\Rightarrow$ 无对流项）、「有效参数」论证（$E_a=32.0$ kJ/mol 介于纯水自扩散 18--20 kJ/mol 与化学吸附之间）。

\textbf{问题 16}：8.1「模型优点」只列了四条常规内容，未把上述三点作为「本文特色」突出；摘要也把最独特的守恒性检验埋在了长段落里。建议在摘要末句和 8.1 中把「有效参数观点 + 五重检验」作为旗帜性表述。

\subsection{4.3 结果的正确性（优，但有三个中等问题）}

四问结果齐全，表格与文件输出完全合规（四位小数、模板格式、「---」与留空处理、result4 的「药材表面」列）；$t^*=56.97$ h、$50.69$ h 均落于题述「2--3 天」区间，并给了问题 3 的不确定度（$\pm0.1$ h）。五重检验属于超规格配置。

\textbf{问题 4（空间精度）}：7.2 节报告 $\Delta\zeta=1/20\to1/40$ 最大差 $5.33\times10^{-4}$，按二阶格式 Richardson 外推，细网格误差 $E_{\mathrm{fine}}\approx5.33\times10^{-4}/3\approx1.8\times10^{-4}$，\textbf{大于}四位小数分辨率 $10^{-4}$。即表 6 中 6 h 行（含水率梯度最陡、收缩最快，是全场最难算的时刻）的第 4 位小数可能偏差 1--2 个单位。12 h 之后误差迅速降至 $10^{-5}$ 量级，四位小数可靠；对 $t^*$ 与物理结论几乎无影响。建议补做 $\Delta\zeta=1/80$（$N=81$）复验：若 1/40 与 1/80 之差降至 $\sim2\times10^{-5}$（约 $5.33\times10^{-4}$ 的 $1/9$），即可用数据证明细网格已进入 $10^{-5}$ 量级；若降不下去，说明另有主导误差源（时间步或表面处理），需排查。

\textbf{问题 5（$t^*$ 收敛）}：问题 3 的验证报了 $t^*$ 的网格收敛（$56.7069$ h @ 0.5 mm vs $56.9674$ h @ 0.25 mm），问题 4 未报。$t^*$ 是问题 4 的主结论，应同样做 $\Delta\zeta=1/20$ 与 $1/40$（或 $1/40$ 与 $1/80$）的对比，确认 $50.6931$ h 稳定。

\textbf{问题 6（归因混杂）}：6.4 节的对比是「附录 3 物性+无收缩」对「附录 4 物性+有收缩」，同时改变了物性（附录 4 的 $D$ 小 2--5 倍，使干燥变慢）与几何收缩（$R:2\to1.198$ cm，使干燥变快）两个因素，6.3 h 的净差不能单独归因于收缩。表 5/表 6 中心含水率曲线在约 42 h 处交叉（前 40 h 问题 4 反而更慢、之后反超），恰好说明两种效应在竞争。建议补一组消融实验：「附录 4 物性 + 固定 $R=2$ cm」得 $t^*_{\mathrm{无收缩}}$，则 $\Delta t_{\mathrm{收缩}}=t^*_{\mathrm{无收缩}}-50.6931$ h 可干净地分离出收缩的单独贡献，使 6.4 节的定性结论变为定量结论。

\subsection{4.4 表述的清晰程度（中，主要失分点）}

\textbf{问题 3（摘要三改）}：其一，摘要已写到页满（36 行），再改即溢出——规范要求摘要原则上不超过一页，必须留冗余；其二，摘要没有问题 2 的结果——问题 1 给了 $36.79\,^\circ$C/1.51 kg/kg，问题 3、4 给了 $t^*$，唯独问题 2 只写「共享同一模型」，评委看摘要应看到每一问的关键答案，建议补一句（如「3 h 时表面降至 1.01 kg/kg、中心 1.77 kg/kg」）；其三，「数值方法与检验」段建议压缩约 40\%，保留 Bessel 四位小数一致、收敛比 4、守恒 $10^{-18}$、$E_a=32.0$ 与文献一致四个要点即可。

\textbf{问题 8（7.6 措辞）}：「若显式计入潜热，时长 79 h 超出 2--3 天上限，恰从反面印证……」——题目说「一般持续 2--3 天」是背景性描述而非硬约束，评委可能认为这是循环论证。建议改为中性表述：把 $+40\%\sim50\%$ 作为「模型不确定度区间」保留在结论中，不作为否定显式潜热模型的证据。

\textbf{问题 13（尺度表述）}：5.5 节给边界层 $\delta\sim D/h_m\approx5$ mm，6.1 节给扩散深度 $\sqrt{Dt}\approx2.6$ mm，评委可能疑惑为何「水分集中于约 5 mm」而扩散深度只有 2.6 mm。建议 6.1 节补一句两者的关系（$\sqrt{Dt}$ 是扩散前沿尺度，$\delta$ 是表面 Robin 边界层厚度，同量级，共同刻画「变化集中于表层」）。

\textbf{问题 15（结论清单）}：建议文末加「本文主要结论」清单——四个问题的答案各一行（含 $t^*$ 与不确定度），评委翻到最后可直接核对。

\textbf{问题 7（公式记号）}：式 (cn) 左端乘 $v_i$，而式 (L) 的 $(Lu)_i$ 已除以 $v_i$，量纲上多出一个 $v_i$。附录 3 的守恒推导（式 (cn) 两边同乘 $v_i$ 求和得式 (conserv)）与代码（散度已除体积，左端容量为 $\rho c_p$）均按「式 (cn) 无 $v_i$」的版本自洽，说明正文公式只是记号残留。二选一即可：去掉式 (cn) 左端的 $v_i$；或去掉式 (L) 末尾的 $/v_i$ 并相应调整圆心行、表面行写法。

\textbf{问题 14}：5.4 节「循环两轮即达收敛」缺量化。建议补一句两轮后耦合残差与时间步误差之比（如 $<1\%$），与 7.2 节的时间收敛数据呼应。

\textbf{问题 11}：表 6 的「到药材中心的距离 $d$」按当前几何位置（欧拉视角）理解，对应 $\zeta=d/R(t)$；题面把「药材表面」单列且要求 $d>R(t)$ 记「---」支持该解读，论文 5.3 节也已写明 $d>R(t)$ 时记「---」，但未明示 $\zeta=d/R(t)$ 的映射关系。补一句显式声明即可消除歧义。

\textbf{问题 12}：result3/4 的末行为 $t^*$ 时刻（$t^*$ 不一定是 60 的整数倍，末行与 60 s 间隔序列不对齐）。建议在附录 5 补一句「末行为烘干结束时刻 $t^*$」。

\section{五、格式规范对照（2026 年修订稿）}

\begin{center}
{\zihao{5}\begin{tabular}{p{6.2cm}p{5.2cm}c}
\toprule
规范要求 & 论文现状 & 结论 \\
\midrule
A4，四边边距 $\ge2.5$ cm & geometry 2.5 cm & 符合 \\
摘要页：题目+摘要+关键词，$\le1$ 页 & 第 1 页，恰好一页 & 已到页满，须留冗余 \\
页码从摘要页起、页脚中部、阿拉伯数字 & 第 1 页底有页码「1」 & 符合 \\
论文不能有页眉，不要目录 & 无页眉无目录 & 符合 \\
正文（不含附录）$\le30$ 页 & 正文约 16 页 & 符合 \\
电子版第一页必须是摘要页 & 第 1 页即摘要页 & 符合 \\
通行惯例：题目三号黑体；一级标题四号黑体居中；二三级小四黑体左对齐；正文小四宋体 & 题目二号黑体；一级标题三号黑体左对齐 & 建议改 \\
参考文献 GB/T 7714 & 18 篇，格式规范 & 符合 \\
\bottomrule
\end{tabular}}
\end{center}

字体问题（\textbf{问题 10}）：官方规范对字号不强制，但通行惯例如此，按惯例改最稳妥。

\section{六、数据核查记录}

\begin{center}
{\zihao{5}\begin{tabular}{p{6.6cm}ll}
\toprule
核查点 & 论文声明 & 实测核对 \\
\midrule
附件 1 点数与范围 & 241 点，0--14400 s & 242 行（含表头），0--14400 s，60 s 间隔 \\
附件 1 末值 & $T_a=50.165$、$C_a=0.04986$ & 一致 \\
附件 1 在 1800 s & 烘房温度 41.51 $^\circ$C & 41.513 $^\circ$C，0.03307 \\
附件 2 点数与收缩 & 145 点；2$\to$1.198 cm；67 h 起恒定 & 一致；30 min 处 1.873（收缩 6.35\%） \\
表 1 在 1800 s（中心/表面） & 33.5753 / 36.7856 & result1 温度表一致 \\
表 2 在 1800 s（中心/表面） & 2.5500 / 1.5103 & result1 水分表一致 \\
表 3/表 4 在 3 h & 49.8495 / 49.9664；1.7662 / 1.0081 & result2 一致 \\
表 5 末行（中心/表面） & 0.1500 / 0.0525 & result3 末行一致 \\
表 6 末行（中心/表面） & 0.1500 / 0.0525 & result4 末行一致 \\
表 6 中 $d>R(t)$ 留空 & 6 h 时 $R=1.374$ cm，1.5/2 cm 列「---」 & result4 对应单元格为空 \\
result 文件结构 & 1 s/60 s 间隔、0.1 cm 间隔、模板格式 & 行数 1801/10801/3420/3043，列数 22/22/22/23 \\
7.2 节外推 & $t^*\approx57.05$ h，不确定度 0.1 h & $(4\times56.9674-56.7069)/3=57.054$ 吻合 \\
7.6 节 Arrhenius 比值 & $e^{3850(1/323.3-1/314.7)}\approx0.72$ & $=0.721$ 吻合 \\
6.4 节「$D_4$ 比 $D_3$ 小 2--5 倍」 & $D_4/D_3=0.175\,e^{0.15/C}$ & $C=2.55$ 时 5.4 倍、$C=0.15$ 时 2.1 倍，吻合 \\
\bottomrule
\end{tabular}}
\end{center}

\section{七、修改优先级}

\textbf{高优先级}
\begin{enumerate}[label=(\arabic*),leftmargin=3em]
\item 补 7.7 节\textbf{灵敏度分析}（问题 2）。重点量化恒温段假设：按 7.6 节自身的 Arrhenius 标度，$T_a$ 每变化 $1\,^\circ$C，$t^*$ 约变化 4--5\%，恒温段 $T_a$ 的假设误差即使只有 1--2 $^\circ$C，对 $t^*$ 的影响也远大于 $\pm0.1$ h 的离散不确定度——这是全文最需要面对的假设。建议给 $T_a\pm2\,^\circ$C、$C_a\pm20\%$、$h,h_m\pm50\%$ 对 $t^*$ 的影响表；
\item \textbf{摘要三改}（问题 3）：压缩检验段、补问题 2 结果、留出页空；
\item 配图 6--7 张（问题 1，\textbf{暂缓}，待配图方案确定后另行处理）。
\end{enumerate}

\textbf{中优先级}
\begin{enumerate}[label=(\arabic*),leftmargin=3em]
\item 问题 4 空间精度复验（$\Delta\zeta=1/80$）与 $t^*$ 网格收敛补报（问题 4、5）；
\item 消融实验分离收缩贡献（问题 6）；
\item 修正式 (cn) 的 $v_i$ 记号（问题 7）；
\item 7.6 节措辞中性化（问题 8）；
\item 5.2 节声明 $h,h_m$ 沿用附录 2（问题 9）。
\end{enumerate}

\textbf{低优先级}
\begin{enumerate}[label=(\arabic*),leftmargin=3em]
\item 字体惯例调整（问题 10）；8.1 突出特色（问题 16）；
\item 表 6 欧拉视角声明（问题 11）；result 末行说明（问题 12）；
\item 尺度表述统一（问题 13）；Picard 残差量化（问题 14）；结论清单（问题 15）。
\end{enumerate}

\textbf{代码仓库问题（转引自《问题四分析\_发现的问题.md》）}：\texttt{debug\_q4.py:59} 表面通量系数多乘 $1/\Delta\zeta$（相差 40 倍），与正式求解器不一致，虽然不影响正式结果，但留在仓库中极易被误用。建议删除或修正，或改为直接复用 \texttt{q4\_solver.py} 的函数。

\section{八、结论}

论文的模型框架、数值实现与主结论（$t^*=56.97$ h / $50.69$ h）可信，附件数据引用与结果文件全部准确。当前最优先的三项改进是：补灵敏度分析、修摘要、补问题 4 的验证完备性（空间精度复验与 $t^*$ 收敛）；配图按团队安排暂缓。完成上述修改后，本文在「结果的正确性」「假设的合理性」两项上具备冲击更高奖项的实力，表述与格式项可随修改同步补齐。

\medskip
{\zihao{5}\textbf{评审依据来源}：《全国大学生数学建模竞赛论文格式规范（2026 年修订稿）》\url{https://dxs.moe.gov.cn/zx/a/hd_sxjm_gsyw/260702/2046411.shtml}；竞赛官网 \url{https://www.mcm.edu.cn/html_cn/node/4cd596519c9eb9fbd866398f6df0caa3.html}；评审标准整理 \url{https://github.com/landawang-star/shuxue-jianmo-lunwen-pingshen/blob/main/references/%E5%9B%BD%E8%B5%9B%E8%AF%84%E5%AE%A1%E6%A0%87%E5%87%86.md}}

\end{document}
